\documentclass{article}%
\usepackage[T1]{fontenc}%
\usepackage[utf8]{inputenc}%
\usepackage{lmodern}%
\usepackage{textcomp}%
\usepackage{lastpage}%
\usepackage{geometry}%
\geometry{margin=2cm}%
\usepackage{expex}%
\usepackage[T1]{fontenc}%
\usepackage[utf8]{inputenc}%
\usepackage[spanish]{babel}%
%
\title{Análisis Lingüístico: Gramatica Normativa Kaqchikel}%
\author{Generado por el TFM de Elías Carrasco}%
%
\begin{document}%
\normalsize%
\maketitle%
\section{Page 170}%
\label{sec:Page170}%
plural) del objeto y sujeto. La palabra principal incluye primeramente el tiempo/aspecto %
o modo (excepto en el perfectivo que es la última), seguido por el juego absolutivo, %
juego ergativo, la raíz o base en calidad de obligatorios y sufijos de categoría en casos %
específicos (si el verbo es derivado). Como elementos opcionales incluye el movimiento, %
su posición es siempre antes del juego ergativo. Ejemplos: %
\subsection{Marca de tiempo/aspecto/modo:}%
\label{subsec:Marcadetiempo/aspecto/modo}%

%
\par\medskip%
\ex
\textit{Todavía fui a ver al niño.} \\
\begingl
\gla Xb'entz'eta' na kan ri ak'wal. //
\glb X-∅-b'e-n-tz'et-a' na kan ri ak'wal //
\glc \textbf{COM-}\textbf{B3s-}\textbf{MOV-}\textbf{A1s-}\textbf{ver-}\textbf{SC} \textbf{PAR} \textbf{DIR} \textbf{DET} \textbf{niño} //
\glft \textbf{Sintaxis:} [Xb'e-n-tz'et-a']\textsubscript{SVT} //
\endgl
\xe
%
\medskip%
\subsection{Marca el objeto en el verbo:}%
\label{subsec:Marcaelobjetoenelverbo}%

%
\par\medskip%
\ex
\textit{...vio muy hermosas a las niñas...} \\
\begingl
\gla ...kan jeb'ël ok xerutz'ët ri ch'uta'q xtani' rija'... //
\glb kan jeb'ël ok x-e-ru-tz'ët ri ch'uta'q xtan-i' rija'... //
\glc \textbf{PAR} \textbf{hermosa} \textbf{DIR} \textbf{COM-}\textbf{B3p-}\textbf{A3s-}\textbf{ver} \textbf{DET} \textbf{DIM} \textbf{Niña-PL} \textbf{él} //
\glft \textbf{Sintaxis:} [xerutz'ët]\textsubscript{SVT} //
\endgl
\xe
%
\medskip%
\subsection{Marca el sujeto en el verbo:}%
\label{subsec:Marcaelsujetoenelverbo}%

%
\par\medskip%
\ex
\textit{Los dos hombres compraron un corte.} \\
\begingl
\gla Xkilöq' jun uqaj ri ka'i' achi'a'. //
\glb X-∅-ki-löq' jun uqaj ri ka'i' achi'-a'. //
\glc \textbf{COM-}\textbf{B3s-}\textbf{A3p-}\textbf{comprar} \textbf{IND} \textbf{corte} \textbf{DET} \textbf{dos} \textbf{hombre-PL} //
\glft \textbf{Sintaxis:} [Xkilöq']\textsubscript{SVT} //
\endgl
\xe
%
\medskip%
\par\medskip%
\ex
\textit{La señorita regaló el tomate.} \\
\begingl
\gla Xusipaj ri xkoya' ri xtän. //
\glb X-∅-u-sip-a-j ri xkoya' ri xtän. //
\glc \textbf{COM-}\textbf{B3s-}\textbf{A3s-}\textbf{regalar-BV-SC} \textbf{DET} \textbf{tomate} \textbf{DET} \textbf{señorita} //
\glft \textbf{Sintaxis:} [Xusipaj]\textsubscript{SVT} //
\endgl
\xe
%
\medskip%
\subsection{Marca movimiento en el verbo:}%
\label{subsec:Marcamovimientoenelverbo}%

%
\par\medskip%
\ex
\textit{...nuestro papá va a trabajar, vamos a ayudarlo,...} \\
\begingl
\gla ...nb'esamäj qatata', nb'eqato',... //
\glb  //
\glc  //
\glft \textbf{Sintaxis:} [nb'eqato']\textsubscript{SVT} //
\endgl
\xe
%
\medskip

%
\end{document}